\documentclass[12pt,a4paper]{article}

\usepackage[T2A]{fontenc}
\usepackage[utf8]{inputenc}
\usepackage[russian]{babel}
\usepackage{amsmath,amssymb,amsthm,mathtools}
\usepackage{geometry}
\usepackage{microtype}
\usepackage{enumitem}
\usepackage{xcolor}
\usepackage{tikz}
\usepackage{hyperref}

\geometry{left=25mm,right=20mm,top=22mm,bottom=24mm}
\hypersetup{
  colorlinks=true,
  linkcolor=blue!45!black,
  urlcolor=blue!55!black
}
\setlist{nosep}
\setlength{\parindent}{1.25em}
\setlength{\parskip}{0.15em}

\newtheorem{theorem}{Теорема}
\newtheorem{proposition}{Утверждение}
\theoremstyle{definition}
\newtheorem{definition}{Определение}
\theoremstyle{remark}
\newtheorem*{remark}{Замечание}

\DeclareMathOperator{\Sp}{Sp}
\DeclareMathOperator{\diag}{diag}
\DeclareMathOperator{\cond}{cond}
\newcommand{\CC}{\mathbb{C}}
\newcommand{\RR}{\mathbb{R}}
\newcommand{\ZZ}{\mathbb{Z}}
\newcommand{\eps}{\varepsilon}
\newcommand{\norm}[1]{\left\lVert #1\right\rVert}

\title{Операторы свёртки, смаз изображения\\
и быстрое преобразование Фурье}
\author{Конспект лекции}
\date{26 сентября 2026 года}

\begin{document}

\maketitle

\section{Рабочее пространство и обозначения}

Продолжаем изучать свёртки на группе $\ZZ_n=\{0,1,\dots,n-1\}$.
Векторы $x=(x_0,\dots,x_{n-1})^{\mathsf T}$ принадлежат $\CC^n$,
а вычитание индексов выполняется по модулю $n$:
\[
  j\ominus r=(j-r)\bmod n.
\]
Через $e_j$ обозначаем стандартный базисный вектор, у которого на месте
$j$ стоит единица, а остальные координаты равны нулю. Напомним формулу
циклической свёртки:
\[
  (a*x)_j=\sum_{r=0}^{n-1}a_{j\ominus r}x_r.
\]
Для фиксированного $a\in\CC^n$ оператор
\[
  A\colon\CC^n\longrightarrow\CC^n,\qquad Ax=a*x,
\]
называется оператором свёртки с ядром $a$. На предыдущей лекции были
доказаны коммутативность и ассоциативность свёртки, а также равенство
$a*e_0=a$.

Матрица циклического сдвига $V$ действует по правилу
\[
  Vx=(x_{n-1},x_0,\dots,x_{n-2})^{\mathsf T},
  \qquad (Vx)_j=x_{j\ominus1}.
\]

\section{Свёртка с базисным вектором и матрица оператора}

\begin{proposition}
Свёртка с вектором $e_1$ совпадает с циклическим сдвигом:
\[
  e_1*x=Vx.
\]
Здесь предполагается $n\geqslant2$.
\end{proposition}

\begin{proof}
По коммутативности свёртки
\[
  (e_1*x)_j=(x*e_1)_j
  =\sum_{r=0}^{n-1}x_{j\ominus r}(e_1)_r.
\]
У вектора $e_1$ ненулевая координата только одна: $(e_1)_1=1$.
Поэтому сумма равна $x_{j\ominus1}=(Vx)_j$.
\end{proof}

\begin{theorem}
Матрица оператора $Ax=a*x$ в стандартном базисе является циклической
матрицей с первым столбцом $a$:
\[
  A=\begin{pmatrix}
    a_0 & a_{n-1} & a_{n-2} & \cdots & a_1\\
    a_1 & a_0 & a_{n-1} & \cdots & a_2\\
    a_2 & a_1 & a_0 & \ddots & \vdots\\
    \vdots & \vdots & \ddots & \ddots & a_{n-1}\\
    a_{n-1} & a_{n-2} & \cdots & a_1 & a_0
  \end{pmatrix}.
\]
\end{theorem}

\begin{proof}
Столбец с номером $j$ --- это вектор $Ae_j$. Для первого столбца
\[
  Ae_0=a*e_0=a.
\]
При $j=0,\dots,n-2$ имеем $e_{j+1}=Ve_j=e_1*e_j$. Используя
ассоциативность, коммутативность и доказанную лемму, получаем
\begin{align*}
  Ae_{j+1}&=a*e_{j+1}=a*(e_1*e_j)\\
           &=e_1*(a*e_j)=V(Ae_j).
\end{align*}
Значит, каждый следующий столбец получается циклическим сдвигом
предыдущего. Тот же вывод следует непосредственно из формулы свёртки:
$A_{jr}=a_{j\ominus r}$.
\end{proof}

Символ оператора свёртки совпадает с символом его циклической матрицы:
\[
  a(t)=a_0+a_1t+\dots+a_{n-1}t^{n-1},\qquad A=a(V).
\]

\section{Матрица дискретного преобразования Фурье}
\label{sec:fourier}

Чтобы обосновать формулы для нормы и числа обусловленности, сначала
разберём свойства матрицы Фурье. В рукописи этот блок расположен после
примера со смазом изображения.

Положим
\[
  \eps_k=e^{2\pi i k/n},\qquad
  v_k=(1,\eps_k^{-1},\dots,\eps_k^{-(n-1)})^{\mathsf T}.
\]
Матрица Фурье состоит из столбцов $v_k$:
\[
  F=(v_0\ v_1\ \cdots\ v_{n-1}),\qquad
  F_{jk}=\eps_k^{-j}=e^{-2\pi i jk/n}.
\]
Знак в показателе отрицательный; нормирующего множителя в $F$ нет.
Это то же соглашение, которое использовалось в предыдущем конспекте.

\subsection{Симметричность}

\begin{proposition}
Матрица $F$ симметрична: $F^{\mathsf T}=F$.
\end{proposition}

\begin{proof}
Для любых $j,k$ имеем
\[
  F_{jk}=e^{-2\pi i jk/n}=e^{-2\pi i kj/n}=F_{kj}.
\]
\end{proof}

Симметричность здесь означает транспонирование без комплексного
сопряжения. В общем случае $F$ не является эрмитовой матрицей.

\subsection{Ортогональность столбцов и строк}

\begin{proposition}
Столбцы матрицы $F$ попарно ортогональны, и норма каждого столбца равна
$\sqrt n$. Те же свойства выполнены для строк. В частности,
\[
  \boxed{F^*F=nE}.
\]
\end{proposition}

\begin{proof}
Используем скалярное произведение
$\langle x,y\rangle=\sum_r\overline{x_r}y_r$. Тогда
\[
  \langle v_k,v_m\rangle
  =\sum_{r=0}^{n-1}\overline{\eps_k^{-r}}\eps_m^{-r}
  =\sum_{r=0}^{n-1}\left(e^{2\pi i(k-m)/n}\right)^r.
\]
Если $k=m$, каждое слагаемое равно единице, и сумма равна $n$.
Если $k\ne m$, положим $q=e^{2\pi i(k-m)/n}$. Тогда $q\ne1$, $q^n=1$,
поэтому сумма геометрической прогрессии равна
\[
  1+q+\dots+q^{n-1}=\frac{q^n-1}{q-1}=0.
\]
Следовательно,
\[
  \langle v_k,v_m\rangle=
  \begin{cases}n,&k=m,\\0,&k\ne m.\end{cases}
\]
Эти скалярные произведения являются элементами $F^*F$.
Свойство строк следует из симметричности $F$.
\end{proof}

\subsection{Обратная матрица и унитарная нормировка}

Из $F^*F=nE$ следует
\[
  \boxed{F^{-1}=\frac1nF^*=\frac1n\overline F}.
\]
Последнее равенство использует симметричность: $F^*=\overline F^{\mathsf T}
=\overline F$. Таким образом, для обращения матрицы Фурье достаточно
изменить знак показателя экспоненты и разделить результат на $n$.

Нормированная матрица
\[
  U=\frac1{\sqrt n}F
\]
унитарна:
\[
  U^*U=UU^*=E.
\]
Поэтому она сохраняет евклидову норму:
\[
  \norm{Ux}_2^2=x^*U^*Ux=x^*x=\norm{x}_2^2.
\]
Для ненормированной матрицы имеем
$\norm{Fx}_2=\sqrt n\,\norm{x}_2$.

\begin{proposition}
Умножение матрицы слева или справа на унитарную матрицу не меняет
подчинённую евклидову норму. В частности,
\[
  \norm{UB}_2=\norm{B}_2,\qquad
  \norm{BU}_2=\norm{B}_2,\qquad
  \norm{UBU^*}_2=\norm{B}_2.
\]
\end{proposition}

\begin{proof}
Первое равенство следует из $\norm{UBx}_2=\norm{Bx}_2$.
Для второго заметим, что $U$ взаимно однозначно переводит единичную
сферу в себя. Поэтому
\[
  \sup_{\norm{x}_2=1}\norm{BUx}_2
  =\sup_{\norm{y}_2=1}\norm{By}_2.
\]
Последнее равенство получается применением первых двух.
\end{proof}

\section{Нормы и число обусловленности оператора свёртки}

\subsection{Нормы по строкам и столбцам}

\begin{theorem}
Для циклической матрицы $A$ с первым столбцом $a$
\[
  \boxed{\norm{A}_1=\norm{A}_\infty=\sum_{r=0}^{n-1}|a_r|}.
\]
\end{theorem}

\begin{proof}
Напомним общие формулы
\[
  \norm{A}_1=\max_j\sum_r|A_{rj}|,\qquad
  \norm{A}_\infty=\max_j\sum_r|A_{jr}|.
\]
В циклической матрице каждый столбец является перестановкой первого
столбца, а каждая строка содержит те же элементы $a_0,\dots,a_{n-1}$.
Следовательно, все суммы модулей по строкам и столбцам одинаковы.
\end{proof}

\subsection{Евклидова норма}

\begin{theorem}
Если $A=a(V)$, то
\[
  \boxed{\norm{A}_2=\max_{0\leqslant k<n}|a(\eps_k)|}.
\]
\end{theorem}

\begin{proof}
По результатам предыдущей лекции
\[
  A=FDF^{-1},\qquad
  D=\diag\bigl(a(\eps_0),\dots,a(\eps_{n-1})\bigr).
\]
Поскольку $F=\sqrt n\,U$ и $F^{-1}=U^*/\sqrt n$, множители
$\sqrt n$ сокращаются:
\[
  A=UDU^*.
\]
По унитарной инвариантности нормы $\norm{A}_2=\norm{D}_2$.
Для диагональной матрицы
\[
  \norm{Dx}_2^2=\sum_k|a(\eps_k)|^2|x_k|^2
  \leqslant\left(\max_k|a(\eps_k)|^2\right)\norm{x}_2^2.
\]
Равенство достигается на базисном векторе, отвечающем максимальному
диагональному элементу по модулю. Теорема доказана.
\end{proof}

\begin{remark}
Произвольное преобразование подобия $B\mapsto PBP^{-1}$ не обязано
сохранять норму. В данном доказательстве существенно, что матрицу Фурье
можно нормировать до унитарной.
\end{remark}

\subsection{Число обусловленности}

\begin{theorem}
Если циклическая матрица $A$ обратима, то её число обусловленности в
евклидовой норме равно
\[
  \boxed{\cond_2(A)=\norm{A}_2\norm{A^{-1}}_2
  =\frac{\max_k|a(\eps_k)|}{\min_k|a(\eps_k)|}}.
\]
\end{theorem}

\begin{proof}
Обратимость равносильна условиям $a(\eps_k)\ne0$ для всех $k$.
Из $A=UDU^*$ получаем $A^{-1}=UD^{-1}U^*$, откуда
\[
  \norm{A^{-1}}_2
  =\max_k\frac1{|a(\eps_k)|}
  =\frac1{\min_k|a(\eps_k)|}.
\]
Умножая на найденное значение $\norm{A}_2$, получаем формулу.
\end{proof}

Чем меньше минимальный модуль собственного числа по сравнению с
максимальным, тем больше число обусловленности. Поэтому сама по себе
обратимость ещё не гарантирует устойчивость восстановления данных.

\section{Применение: восстановление смазанного изображения}

\subsection{Модель горизонтального смаза}

Рассмотрим одну строку изображения. Число $x_j$ задаёт яркость пикселя
с номером $j$. При равномерном смещении изображения во время экспозиции
наблюдаемая яркость может моделироваться средним нескольких соседних
значений. Для смаза на два пикселя:
\[
  y_j=\frac{x_j+x_{j-1}}2.
\]
Например,
\[
  y_0=\frac{x_0+x_{-1}}2,\qquad
  y_1=\frac{x_1+x_0}2.
\]
На этом этапе отрицательные индексы обозначают соседние пиксели слева.
Для конечной строки далее требуется выбрать правило обработки границ.

\begin{definition}
В рассматриваемой модели \emph{смазом на $k$ пикселей} называется
усреднение
\[
  y_j=\frac1k\bigl(x_j+x_{j-1}+\dots+x_{j-k+1}\bigr).
\]
\end{definition}

Восстановление изображения означает решение обратной задачи: по
известным значениям $y_j$ найти исходные $x_j$.

Усреднение может терять информацию. В частности, две разные
двусторонние последовательности
\[
  x_j=\frac12\quad\text{для всех }j,
  \qquad
  \widetilde x_j=\begin{cases}0,&j\text{ чётно},\\1,&j\text{ нечётно}\end{cases}
\]
после смаза на два пикселя дают одну и ту же постоянную последовательность
$y_j=1/2$. Для конечной циклической строки этот пример работает при
чётном числе пикселей.

\subsection{Циклическая модель и матрица смаза}

Для круговой панорамы естественно считать, что за последним пикселем
снова идёт первый. Поэтому далее используем циклические индексы:
\[
  y_j=\frac1k\sum_{r=0}^{k-1}x_{j\ominus r},
  \qquad j=0,\dots,n-1,\qquad 1\leqslant k\leqslant n.
\]
Это дополнительное условие модели: произвольная конечная строка
изображения сама по себе не обязана быть периодической.

\begin{figure}[ht]
  \centering
  \begin{tikzpicture}[scale=1.2]
    \draw[gray!70] (0,0) circle (1.15);
    \foreach \j/\ang in {0/90,1/30,2/-30,3/-90,4/-150,5/150} {
      \fill[blue!60!black] (\ang:1.15) circle (0.045);
      \node at (\ang:1.52) {$x_{\j}$};
    }
    \draw[->,thick,blue!60!black] (145:0.85)
      arc[start angle=145,end angle=95,radius=0.85];
    \node[align=left] at (4.1,0.2)
      {$n=6,\quad k=3$\\[5pt]
       $y_0=(x_0+x_5+x_4)/3$\\[5pt]
       $y_1=(x_1+x_0+x_5)/3$};
  \end{tikzpicture}
  \caption{Циклическое усреднение: границы строки соединены}
\end{figure}

Положим
\[
  s_r=\begin{cases}1/k,&0\leqslant r<k,\\0,&k\leqslant r<n.\end{cases}
\]
Тогда $y=s*x=Sx$, где
\[
  S=\frac1k(E+V+\dots+V^{k-1}).
\]
Первый столбец $S$ содержит $k$ чисел $1/k$, за которыми следуют нули.
Символ матрицы смаза равен
\[
  \boxed{s(t)=\frac1k(1+t+\dots+t^{k-1})}.
\]
Таким образом, задача восстановления принимает вид $Sx=y$.

\subsection{Критерий обратимости смаза}

\begin{theorem}
Матрица циклического смаза на $k$ пикселей для строки длины $n$ обратима
тогда и только тогда, когда $n$ и $k$ взаимно просты:
\[
  \boxed{S\text{ обратима}\quad\Longleftrightarrow\quad\gcd(n,k)=1}.
\]
\end{theorem}

\begin{proof}
По критерию обратимости циклической матрицы нужно проверить значения
$s(\eps_\ell)$, $\ell=0,\dots,n-1$. При $\ell=0$ имеем
$\eps_0=1$ и $s(1)=1$. Для $\ell\ne0$ используем сумму геометрической
прогрессии:
\[
  s(\eps_\ell)=\frac1k\frac{1-\eps_\ell^k}{1-\eps_\ell}.
\]
Знаменатель ненулевой. Поэтому
\[
  s(\eps_\ell)=0\quad\Longleftrightarrow\quad\eps_\ell^k=1.
\]
Нулевое собственное число существует ровно тогда, когда у уравнений
$z^n=1$ и $z^k=1$ есть общий корень $z\ne1$.

Если $d=\gcd(n,k)>1$, запишем $n=dn_1$, $k=dk_1$ и возьмём
\[
  z=e^{2\pi i/d}=\eps_{n_1}.
\]
Тогда $z\ne1$, но $z^n=z^k=1$. Следовательно, $S$ необратима.

Обратно, если $\gcd(n,k)=1$, существуют целые $p,q$ такие, что
$pn+qk=1$. Для любого общего корня $z^n=z^k=1$ получаем
\[
  z=z^{pn+qk}=(z^n)^p(z^k)^q=1.
\]
Других общих корней нет. Значит, ни одно из значений $s(\eps_\ell)$
не равно нулю, и матрица $S$ обратима.
\end{proof}

В частности, при $k=1$ имеем $S=E$. При $k=2$ восстановление в
циклической модели однозначно для нечётного $n$ и неоднозначно для
чётного $n$, если решение существует. При $k=n>1$ все координаты
результата равны среднему исходных значений, поэтому восстановить
произвольный исходный вектор невозможно.

\begin{remark}
Для вырожденной матрицы $S$ система $Sx=y$ может быть несовместной.
Если решение есть, оно не единственно: к нему можно прибавить любой
вектор из ядра $S$. Когда $y$ получен точным смазом некоторого $x$,
существование решения обеспечено самим построением.
\end{remark}

\section{Преобразование Фурье и теорема о свёртке}

\subsection{Прямое и обратное преобразования}

Дискретным преобразованием Фурье вектора $x$ называем
\[
  \widehat x=Fx,\qquad
  \widehat x_k=\sum_{r=0}^{n-1}x_r\eps_k^{-r}.
\]
Обратное преобразование имеет вид
\[
  x=F^{-1}\widehat x=\frac1nF^*\widehat x,\qquad
  x_r=\frac1n\sum_{k=0}^{n-1}\widehat x_k\eps_k^r.
\]
Для двух векторов $u,v$ обозначим через $u\odot v$ их покоординатное
произведение: $(u\odot v)_k=u_kv_k$.

\begin{theorem}[о свёртке]
Для любых $a,b\in\CC^n$ выполнено
\[
  \boxed{F(a*b)=(Fa)\odot(Fb)}.
\]
То есть преобразование Фурье переводит циклическую свёртку в
покоординатное умножение.
\end{theorem}

\begin{proof}
Для координаты с номером $k$ запишем
\begin{align*}
  \bigl(F(a*b)\bigr)_k
  &=\sum_{r=0}^{n-1}(a*b)_r\eps_k^{-r}\\
  &=\sum_{r=0}^{n-1}\sum_{m=0}^{n-1}
     a_{r\ominus m}b_m\eps_k^{-r}\\
  &=\sum_{m=0}^{n-1}b_m
     \sum_{r=0}^{n-1}a_{r\ominus m}\eps_k^{-r}.
\end{align*}
При фиксированном $m$ сделаем замену $j=r\ominus m$, то есть
$r\equiv j+m\pmod n$. Поскольку $\eps_k^n=1$,
\[
  \eps_k^{-r}=\eps_k^{-(j+m)}=\eps_k^{-j}\eps_k^{-m}.
\]
Следовательно,
\begin{align*}
  \bigl(F(a*b)\bigr)_k
  &=\sum_{m=0}^{n-1}b_m\eps_k^{-m}
    \sum_{j=0}^{n-1}a_j\eps_k^{-j}\\
  &=(Fb)_k(Fa)_k.
\end{align*}
Равенство выполнено для всех $k$.
\end{proof}

Отсюда получается способ вычисления свёртки:
\[
  a*b=F^{-1}\bigl((Fa)\odot(Fb)\bigr).
\]
В принятых соглашениях в прямой формуле теоремы нет дополнительного
множителя; коэффициент $1/n$ входит в обратное преобразование.

\begin{remark}
Не следует смешивать символ $a(t)$ и преобразование $Fa$ с одинаковыми
индексами:
\[
  (Fa)_k=a(\eps_k^{-1}),\qquad
  a(\eps_k)=(Fa)_{(-k)\bmod n}.
\]
Поэтому диагонализация из предыдущей лекции и теорема о свёртке
согласованы, но используют противоположный порядок частот:
\[
  F^{-1}AF=\diag\bigl(a(\eps_k)\bigr),\qquad
  FAF^{-1}=\diag(Fa).
\]
\end{remark}

Для модели смаза $Sx=y$ теорема даёт $\widehat y=\widehat s\odot\widehat x$.
Если $S$ обратима, то
\[
  \widehat x_k=\frac{\widehat y_k}{\widehat s_k},\qquad
  x=F^{-1}\widehat x.
\]
Эта формула поясняет, почему малые значения $|\widehat s_k|$ усиливают
ошибки в данных. Она является непосредственным следствием доказанных
утверждений.

\section{Быстрое преобразование Фурье}

\subsection{Прямое вычисление и разбиение отсчётов}

В этом разделе размер вектора обозначаем $N$, как в рукописи.
Пусть $N=2^m$ и
\[
  \eps=e^{2\pi i/N},\qquad
  X_k=\sum_{r=0}^{N-1}x_r\eps^{-rk},\qquad k=0,\dots,N-1.
\]
Прямое вычисление содержит $N$ сумм по $N$ слагаемых и требует
$O(N^2)$ арифметических операций. Быстрое преобразование Фурье (БПФ)
вычисляет тот же результат, используя повторяющуюся структуру этих сумм.

Положим $h=N/2$ и разделим отсчёты на чётные и нечётные:
\[
  (x_0,x_2,\dots,x_{N-2}),\qquad
  (x_1,x_3,\dots,x_{N-1}).
\]
Тогда
\begin{align*}
  X_k
  &=\sum_{r=0}^{h-1}x_{2r}\eps^{-2rk}
    +\sum_{r=0}^{h-1}x_{2r+1}\eps^{-(2r+1)k}\\
  &=\sum_{r=0}^{h-1}x_{2r}(\eps^2)^{-rk}
    +\eps^{-k}\sum_{r=0}^{h-1}x_{2r+1}(\eps^2)^{-rk}.
\end{align*}
Поскольку
\[
  \eps^2=e^{2\pi i/h},
\]
обе суммы являются преобразованиями Фурье длины $h$. Обозначим их
\[
  E_k=\sum_{r=0}^{h-1}x_{2r}(\eps^2)^{-rk},\qquad
  O_k=\sum_{r=0}^{h-1}x_{2r+1}(\eps^2)^{-rk}.
\]
Итак,
\[
  X_k=E_k+\eps^{-k}O_k.
\]
Буква $E_k$ здесь обозначает преобразование чётных отсчётов, а не
единичную матрицу.

\subsection{Получение второй половины результата}

Последовательности $E_k$ и $O_k$ периодичны по индексу $k$ с периодом
$h$. Например,
\begin{align*}
  E_{k+h}
  &=\sum_{r=0}^{h-1}x_{2r}(\eps^2)^{-r(k+h)}\\
  &=\sum_{r=0}^{h-1}x_{2r}(\eps^2)^{-rk}
    \underbrace{(\eps^2)^{-rh}}_{=1}=E_k.
\end{align*}
Аналогично $O_{k+h}=O_k$. Кроме того, $\eps^h=-1$, поэтому
\[
  \eps^{-(k+h)}=-\eps^{-k}.
\]
Следовательно, для $k=0,\dots,h-1$ получаем две формулы:
\[
  \boxed{\begin{aligned}
    X_k&=E_k+\eps^{-k}O_k,\\
    X_{k+h}&=E_k-\eps^{-k}O_k.
  \end{aligned}}
\]
Одна пара значений $E_k,O_k$ даёт сразу два элемента полного
преобразования длины $N$.

\subsection{Операция <<бабочка>>}

Положим $T_k=\eps^{-k}O_k$. Тогда объединение двух преобразований
половинной длины записывается как
\[
  (E_k,O_k)\longmapsto(E_k+T_k,E_k-T_k).
\]
Для каждой пары требуется одно умножение на коэффициент $\eps^{-k}$,
одно сложение и одно вычитание. Эту операцию называют <<бабочкой>>.

\begin{figure}[ht]
  \centering
  \begin{tikzpicture}[x=1cm,y=1cm,>=stealth]
    \node[left] at (0,1.6) {$E_k$};
    \node[left] at (0,0) {$O_k$};
    \node[draw,rounded corners,fill=blue!5,minimum width=1.1cm,
      minimum height=0.55cm] (mul) at (1.1,0) {$\eps^{-k}$};
    \draw[->,thick] (0,0) -- (mul.west);
    \draw[thick] (mul.east) -- (2.1,0);
    \node[below] at (2.1,0) {$T_k$};
    \draw[thick] (0,1.6) -- (2.1,1.6);
    \fill (2.1,0) circle (0.035);
    \fill (2.1,1.6) circle (0.035);
    \node[draw,circle,inner sep=2pt] (sum) at (4.1,1.6) {$+$};
    \node[draw,circle,inner sep=2pt] (diff) at (4.1,0) {$-$};
    \draw[->,thick,blue!60!black] (2.1,1.6) -- (sum.west);
    \draw[->,thick,blue!60!black] (2.1,0) -- (sum.south west);
    \draw[->,thick,blue!60!black] (2.1,1.6) -- (diff.north west);
    \draw[->,thick,blue!60!black] (2.1,0) -- (diff.west);
    \draw[->,thick] (sum.east) -- (5.1,1.6);
    \draw[->,thick] (diff.east) -- (5.1,0);
    \node[right] at (5.1,1.6) {$X_k=E_k+T_k$};
    \node[right] at (5.1,0) {$X_{k+h}=E_k-T_k$};
  \end{tikzpicture}
  \caption{Объединение результатов чётной и нечётной ветвей БПФ}
\end{figure}

\subsection{Пример: вектор из четырёх отсчётов}

Пусть $x=(1,2,3,4)^{\mathsf T}$. Тогда $N=4$, $h=2$ и $\eps=i$.
Разобьём исходный вектор на чётную и нечётную части:
\[
  x_{\mathrm{even}}=(1,3)^{\mathsf T},\qquad
  x_{\mathrm{odd}}=(2,4)^{\mathsf T}.
\]
Матрица Фурье размера $2$ равна
\[
  F_2=\begin{pmatrix}1&1\\1&-1\end{pmatrix}.
\]
Поэтому преобразование вектора $(a,b)^{\mathsf T}$ имеет вид
$(a+b,a-b)^{\mathsf T}$. Получаем
\[
  E=(4,-2)^{\mathsf T},\qquad O=(6,-2)^{\mathsf T}.
\]
Теперь применим <<бабочку>>:
\begin{center}
\renewcommand{\arraystretch}{1.3}
\begin{tabular}{c|c|c|c|c}
  $k$ & $\eps^{-k}$ & $T_k=\eps^{-k}O_k$
      & $X_k=E_k+T_k$ & $X_{k+2}=E_k-T_k$\\\hline
  $0$ & $1$ & $6$ & $10$ & $-2$\\
  $1$ & $-i$ & $2i$ & $-2+2i$ & $-2-2i$
\end{tabular}
\end{center}
Итоговое преобразование:
\[
  \boxed{X=(10,-2+2i,-2,-2-2i)^{\mathsf T}}.
\]

\section{Рекурсия и сложность БПФ}

Разбиение на чётные и нечётные отсчёты повторяется рекурсивно. Для
$N=2^m$ после $m=\log_2N$ шагов остаются задачи длины $1$, для которых
преобразование совпадает с самим числом.

Например, индексы при $N=8$ разбиваются следующим образом:
\begin{figure}[ht]
  \centering
  \begin{tikzpicture}[x=1cm,y=1cm,>=stealth]
    \node (root) at (0,0) {$0,1,2,3,4,5,6,7$};
    \node (even) at (-2.5,-0.9) {$0,2,4,6$};
    \node (odd) at (2.5,-0.9) {$1,3,5,7$};
    \node (a) at (-3.8,-1.8) {$0,4$};
    \node (b) at (-1.2,-1.8) {$2,6$};
    \node (c) at (1.2,-1.8) {$1,5$};
    \node (d) at (3.8,-1.8) {$3,7$};
    \foreach \from/\to in {root/even,root/odd,even/a,even/b,odd/c,odd/d}
      \draw[->,blue!60!black] (\from) -- (\to);
  \end{tikzpicture}
  \caption{Первые уровни рекурсивного разделения отсчётов}
\end{figure}

На глубине $s$ имеется $2^s$ подзадач длины $N/2^s$. На объединение
одной такой задачи требуется $O(N/2^s)$ операций, поэтому суммарная
работа на одном уровне равна
\[
  2^s\,O\left(\frac{N}{2^s}\right)=O(N).
\]
Уровней $\log_2N$, следовательно, общая сложность составляет
\[
  \boxed{O(N\log_2N)}.
\]
То же самое выражает рекуррентное соотношение
\[
  T(N)=2T(N/2)+O(N),\qquad T(1)=O(1).
\]
В сравнении с прямым вычислением $O(N^2)$ выигрыш возникает за счёт
повторного использования результатов преобразований меньшего размера.

\begin{remark}
В конце рукописи отмечен переход от длины $1000$ к $1024=2^{10}$.
Естественная интерпретация --- дополнение данных нулями до степени
двойки для описанного алгоритма. При этом вычисляется ДПФ дополненного
вектора длины $1024$, а не исходное ДПФ длины $1000$: частотная сетка
меняется. Для циклической свёртки длина также определяет период, поэтому
увеличивать её без учёта задачи нельзя. Это уточнение к краткой пометке
в рукописи.
\end{remark}

\end{document}
